%==================================================================
% GALERI ELEMEN
%
% Berkas ini berisi CONTOH SIAP SALIN untuk seluruh elemen yang lazim
% dibutuhkan pada Tugas Akhir: berbagai jenis tabel, gambar, rumus,
% kode program, pseudocode, dan pola sitasi.
%
% CARA MEMATIKANNYA
% Beri tanda % pada baris %==================================================================
% GALERI ELEMEN
%
% Berkas ini berisi CONTOH SIAP SALIN untuk seluruh elemen yang lazim
% dibutuhkan pada Tugas Akhir: berbagai jenis tabel, gambar, rumus,
% kode program, pseudocode, dan pola sitasi.
%
% CARA MEMATIKANNYA
% Beri tanda % pada baris %==================================================================
% GALERI ELEMEN
%
% Berkas ini berisi CONTOH SIAP SALIN untuk seluruh elemen yang lazim
% dibutuhkan pada Tugas Akhir: berbagai jenis tabel, gambar, rumus,
% kode program, pseudocode, dan pola sitasi.
%
% CARA MEMATIKANNYA
% Beri tanda % pada baris %==================================================================
% GALERI ELEMEN
%
% Berkas ini berisi CONTOH SIAP SALIN untuk seluruh elemen yang lazim
% dibutuhkan pada Tugas Akhir: berbagai jenis tabel, gambar, rumus,
% kode program, pseudocode, dan pola sitasi.
%
% CARA MEMATIKANNYA
% Beri tanda % pada baris \include{zz-galeri-elemen} di laporan.tex
% sebelum naskah dikumpulkan. Berkas ini tidak akan ikut tercetak,
% tetapi tetap tersimpan sebagai rujukan.
%==================================================================

% Galeri disertakan setelah lampiran, sehingga tanpa pengaturan ini
% penomorannya akan ikut melanjutkan huruf lampiran (D, E, dst.).
% Blok berikut membuatnya bernomor "G" agar jelas terpisah.
\clearpage
\setcounter{chapter}{0}
\setcounter{figure}{0}
\setcounter{table}{0}
\setcounter{equation}{0}
\renewcommand{\thechapter}{G}
\renewcommand{\thesection}{G.\arabic{section}}
\renewcommand{\thesubsection}{\thesection.\arabic{subsection}}
\renewcommand{\thefigure}{G.\arabic{figure}}
\renewcommand{\thetable}{G.\arabic{table}}
\renewcommand{\theequation}{G.\arabic{equation}}

\chapter*{GALERI ELEMEN (HAPUS SEBELUM DIKUMPULKAN)}
\addcontentsline{toc}{chapter}{GALERI ELEMEN (HAPUS SEBELUM DIKUMPULKAN)}

Bab ini bukan bagian dari naskah Tugas Akhir. Isinya adalah contoh siap salin
untuk seluruh elemen yang lazim dibutuhkan. Matikan dengan memberi tanda
\texttt{\%} pada baris \texttt{\textbackslash include\{zz-galeri-elemen\}} di
berkas \texttt{laporan.tex}.

% ==================================================================
\section{Tabel}
% ==================================================================

\subsection{Tabel biasa}
Keterangan tabel diletakkan di \textbf{atas} tabel. Setiap tabel wajib diberi
\texttt{\textbackslash label} dan disebut dalam teks, misalnya
\cref{tab:g-biasa}.

\begin{table}[H]
    \centering
    \caption{Contoh tabel biasa}
    \label{tab:g-biasa}
    \begin{tabular}{|c|l|c|}
        \hline
        \textbf{No} & \textbf{Keterangan} & \textbf{Nilai} \\ \hline
        1 & Baris pertama & 3,12 \\ \hline
        2 & Baris kedua   & 3,68 \\ \hline
        3 & Baris ketiga  & 3,54 \\ \hline
    \end{tabular}
\end{table}

\subsection{Tabel dengan garis gaya ilmiah (booktabs)}
Sebagian program studi menghendaki tabel tanpa garis vertikal.

\begin{table}[H]
    \centering
    \caption{Contoh tabel gaya booktabs}
    \label{tab:g-booktabs}
    \begin{tabular}{lcc}
        \toprule
        \textbf{Dimensi} & \textbf{Kepentingan} & \textbf{Kinerja} \\
        \midrule
        \textit{Usability}           & 4,41 & 3,12 \\
        \textit{Information quality} & 4,22 & 3,68 \\
        \bottomrule
    \end{tabular}
\end{table}

\subsection{Tabel bersambung antarhalaman (longtable)}
Gunakan \texttt{longtable} bila tabel lebih panjang daripada satu halaman.
Baris kepala akan diulang otomatis, dan halaman sambungan diberi keterangan
``Lanjutan'' sesuai kelaziman naskah Indonesia.

\begin{longtable}{|c|p{6.5cm}|c|c|}
    \caption{Contoh tabel bersambung antarhalaman}
    \label{tab:g-longtable} \\
    \hline
    \textbf{No} & \textbf{Pernyataan} & \textbf{Kepentingan} & \textbf{Kinerja} \\ \hline
    \endfirsthead

    \multicolumn{4}{l}{\footnotesize\textit{Lanjutan \cref{tab:g-longtable}}} \\
    \hline
    \textbf{No} & \textbf{Pernyataan} & \textbf{Kepentingan} & \textbf{Kinerja} \\ \hline
    \endhead

    \hline
    \multicolumn{4}{r}{\footnotesize\textit{Bersambung ke halaman berikutnya}} \\
    \endfoot

    \hline
    \endlastfoot

    1  & Saya merasa sistem mudah dipelajari & 4,44 & 3,10 \\ \hline
    2  & Interaksi dengan sistem jelas dan mudah dipahami & 4,38 & 3,21 \\ \hline
    3  & Sistem mudah dioperasikan & 4,41 & 3,15 \\ \hline
    4  & Navigasi antarhalaman terasa mudah & 4,61 & 3,03 \\ \hline
    5  & Tampilan sistem sesuai dengan jenis layanannya & 4,02 & 3,44 \\ \hline
    6  & Tampilan sistem menarik & 3,88 & 3,52 \\ \hline
    7  & Sistem menyampaikan kesan kompeten & 4,12 & 3,40 \\ \hline
    8  & Sistem menciptakan pengalaman yang menyenangkan & 3,95 & 3,33 \\ \hline
    9  & Informasi yang disajikan akurat & 4,30 & 3,71 \\ \hline
    10 & Informasi yang disajikan tepercaya & 4,28 & 3,69 \\ \hline
    11 & Informasi selalu diperbarui & 4,19 & 3,58 \\ \hline
    12 & Informasi relevan dengan kebutuhan & 4,25 & 3,66 \\ \hline
    13 & Informasi mudah dipahami & 4,21 & 3,74 \\ \hline
    14 & Informasi disajikan pada tingkat rincian yang tepat & 4,08 & 3,70 \\ \hline
    15 & Informasi disajikan dalam format yang sesuai & 4,15 & 3,68 \\ \hline
    16 & Sistem memberikan rasa aman & 4,09 & 3,61 \\ \hline
    17 & Data pribadi terasa terlindungi & 4,18 & 3,55 \\ \hline
    18 & Sistem memudahkan komunikasi dengan pengelola & 3,97 & 3,42 \\ \hline
    19 & Sistem menciptakan rasa keterhubungan & 3,82 & 3,39 \\ \hline
    20 & Layanan diberikan sesuai yang dijanjikan & 4,11 & 3,58 \\ \hline
\end{longtable}

\subsection{Tabel dengan sel gabung (multirow dan multicolumn)}

\begin{table}[H]
    \centering
    \caption{Contoh tabel dengan sel gabung}
    \label{tab:g-gabung}
    \begin{tabular}{|c|c|c|c|c|}
        \hline
        \multirow{2}{*}{\textbf{Butir}} & \multicolumn{2}{c|}{\textbf{Sebelum}} & \multicolumn{2}{c|}{\textbf{Sesudah}} \\ \cline{2-5}
         & \textbf{Rerata} & \textbf{SB} & \textbf{Rerata} & \textbf{SB} \\ \hline
        B1 & 3,10 & 0,82 & 4,05 & 0,61 \\ \hline
        B2 & 3,22 & 0,77 & 4,12 & 0,58 \\ \hline
    \end{tabular}
\end{table}

\subsection{Tabel berkolom otomatis (tabularx)}
Lebar kolom bertipe \texttt{X} menyesuaikan sendiri agar tabel tepat memenuhi
lebar teks.

\begin{table}[H]
    \centering
    \caption{Contoh tabel dengan lebar kolom otomatis}
    \label{tab:g-tabularx}
    \begin{tabularx}{\textwidth}{|c|X|X|}
        \hline
        \textbf{No} & \textbf{Temuan} & \textbf{Tindak lanjut} \\ \hline
        1 & Menu utama tidak konsisten antarhalaman sehingga pengguna kehilangan orientasi. & Menyeragamkan struktur menu pada seluruh halaman. \\ \hline
        2 & Pesan kesalahan tidak menjelaskan penyebab masalah. & Menulis ulang pesan kesalahan dengan bahasa yang spesifik. \\ \hline
    \end{tabularx}
\end{table}

\subsection{Tabel dengan catatan kaki (threeparttable)}
Gunakan bila tabel memerlukan keterangan simbol atau sumber.

\begin{table}[H]
    \centering
    \begin{threeparttable}
        \caption{Contoh tabel dengan catatan kaki}
        \label{tab:g-catatan}
        \begin{tabular}{|l|c|c|}
            \hline
            \textbf{Dimensi} & \textbf{Rerata} & \textbf{Kategori} \\ \hline
            \textit{Usability}\tnote{a}           & 3,12 & Cukup \\ \hline
            \textit{Information quality}          & 3,68 & Baik  \\ \hline
            \textit{Service interaction}\tnote{b} & 3,54 & Baik  \\ \hline
        \end{tabular}
        \begin{tablenotes}\footnotesize
            \item[a] Terdiri atas 8 butir pernyataan.
            \item[b] Terdiri atas 7 butir pernyataan.
            \item Sumber: Hasil pengolahan data penelitian (2026).
        \end{tablenotes}
    \end{threeparttable}
\end{table}

\subsection{Tabel yang terlalu lebar}
Bila tabel sedikit melebihi lebar teks, perkecil jarak antarkolom dengan
\texttt{\textbackslash tabcolsep} atau bungkus dengan \texttt{adjustbox}.

\begin{table}[H]
    \centering
    \caption{Contoh tabel lebar yang dikecilkan otomatis}
    \label{tab:g-lebar}
    \footnotesize
    \begin{adjustbox}{max width=\textwidth}
    \begin{tabular}{|c|c|c|c|c|c|c|c|c|c|c|}
        \hline
        \textbf{Butir} & \textbf{R1} & \textbf{R2} & \textbf{R3} & \textbf{R4} & \textbf{R5} & \textbf{R6} & \textbf{R7} & \textbf{R8} & \textbf{Rerata} & \textbf{Simpangan Baku} \\ \hline
        B1 & 4 & 5 & 4 & 3 & 4 & 5 & 4 & 4 & 4,13 & 0,64 \\ \hline
        B2 & 3 & 4 & 4 & 4 & 5 & 4 & 3 & 4 & 3,88 & 0,64 \\ \hline
    \end{tabular}
    \end{adjustbox}
\end{table}

\subsection{Tabel lanskap}
Paket \texttt{pdflscape} ikut memutar halaman pada berkas PDF, sehingga pembaca
tidak perlu memiringkan kepala.

\begin{landscape}
\begin{table}[H]
    \centering
    \caption{Contoh tabel pada halaman lanskap}
    \label{tab:g-lanskap}
    \footnotesize
    \begin{tabular}{|c|p{4cm}|p{4cm}|p{4cm}|p{4cm}|}
        \hline
        \textbf{No} & \textbf{Kondisi berjalan} & \textbf{Masalah} & \textbf{Usulan} & \textbf{Dampak yang diharapkan} \\ \hline
        1 & Menu tersebar pada tiga tingkat & Pengguna kesulitan menemukan menu & Menyusun ulang menjadi dua tingkat & Waktu pencarian menu berkurang \\ \hline
        2 & Alur pengisian tidak menampilkan tahapan & Pengguna tidak tahu posisinya & Menambahkan penanda tahapan & Kesalahan pengisian berkurang \\ \hline
    \end{tabular}
\end{table}
\end{landscape}

% ==================================================================
\section{Gambar}
% ==================================================================

\subsection{Gambar tunggal}
Keterangan gambar diletakkan di \textbf{bawah} gambar. Lihat
\cref{fig:g-tunggal}.

\begin{figure}[H]
    \centering
    \includegraphics[width=0.6\linewidth]{contoh-diagram}
    \caption{Contoh gambar tunggal}
    \label{fig:g-tunggal}
\end{figure}

\subsection{Gambar dengan keterangan sumber}
Gunakan perintah \texttt{\textbackslash sumber} bila gambar dikutip dari
sumber lain.

\begin{figure}[H]
    \centering
    \includegraphics[width=0.55\linewidth]{contoh-diagram}
    \caption{Contoh gambar yang dikutip dari sumber lain}
    \label{fig:g-sumber}
\end{figure}
\sumber{\citet{barnes2002}}

\subsection{Dua gambar berdampingan}
Gunakan \texttt{subcaption} agar tiap bagian memperoleh label (a) dan (b) yang
dapat dirujuk terpisah, misalnya \cref{fig:g-kiri} dan \cref{fig:g-kanan}.

\begin{figure}[H]
    \centering
    \begin{subfigure}[t]{0.47\linewidth}
        \centering
        \includegraphics[width=\linewidth]{contoh-tangkapan-a}
        \caption{Tampilan sebelum perbaikan}
        \label{fig:g-kiri}
    \end{subfigure}
    \hfill
    \begin{subfigure}[t]{0.47\linewidth}
        \centering
        \includegraphics[width=\linewidth]{contoh-tangkapan-b}
        \caption{Tampilan sesudah perbaikan}
        \label{fig:g-kanan}
    \end{subfigure}
    \caption{Perbandingan tampilan antarmuka}
    \label{fig:g-berdampingan}
\end{figure}

\subsection{Gambar tinggi satu halaman}
Untuk gambar yang tinggi, seperti diagram alir panjang, gunakan
\texttt{height} agar tidak meluber ke halaman berikutnya.

\begin{figure}[H]
    \centering
    \includegraphics[height=0.78\textheight,keepaspectratio]{contoh-tinggi}
    \caption{Contoh gambar tinggi satu halaman}
    \label{fig:g-tinggi}
\end{figure}

\subsection{Gambar lanskap}

\begin{landscape}
\begin{figure}[H]
    \centering
    \includegraphics[width=0.92\linewidth]{contoh-lanskap}
    \caption{Contoh gambar pada halaman lanskap}
    \label{fig:g-lanskap}
\end{figure}
\end{landscape}

\subsection{Catatan penempatan gambar}
Parameter \texttt{[H]} memaksa gambar muncul tepat di posisi penulisannya.
Parameter \texttt{[htbp]} memberi LaTeX kebebasan menempatkan gambar pada
posisi terbaik, sehingga halaman tidak menyisakan ruang kosong. Gunakan
\texttt{[H]} bila urutan gambar penting, dan \texttt{[htbp]} bila
mengutamakan kerapian halaman.

% ==================================================================
\section{Persamaan}
% ==================================================================

Persamaan dalam baris teks ditulis seperti $\bar{x} = \sum x_i / n$, sedangkan
persamaan bernomor ditulis terpisah seperti \cref{eq:g-rerata}.

\begin{equation}
    \label{eq:g-rerata}
    \bar{x} = \frac{1}{n}\sum_{i=1}^{n} x_{i}
\end{equation}

Persamaan bertingkat ditulis dengan lingkungan \texttt{align}, dan setiap baris
memperoleh nomor tersendiri seperti \cref{eq:g-gap} dan \cref{eq:g-tki}.

\begin{align}
    \text{Gap}_i &= \bar{P}_i - \bar{I}_i \label{eq:g-gap} \\
    \text{TKI}_i &= \frac{\bar{P}_i}{\bar{I}_i} \times 100\% \label{eq:g-tki}
\end{align}

Beberapa bentuk yang sering dibutuhkan: pecahan $\frac{a}{b}$, akar
$\sqrt[3]{x}$, jumlah $\sum_{i=1}^{n}$, hasil kali $\prod_{i=1}^{n}$, integral
$\int_{0}^{1}$, huruf Yunani $\alpha, \beta, \sigma, \mu, \kappa$, serta matriks
berikut.

\begin{equation}
    \label{eq:g-matriks}
    \mathbf{D} =
    \begin{bmatrix}
        d_{11} & d_{12} & \cdots & d_{1n} \\
        d_{21} & d_{22} & \cdots & d_{2n} \\
        \vdots & \vdots & \ddots & \vdots \\
        d_{m1} & d_{m2} & \cdots & d_{mn}
    \end{bmatrix}
\end{equation}

% ==================================================================
\section{Kode Program}
% ==================================================================

\subsection{Kode ditulis langsung}
Kode dapat ditulis langsung di dalam naskah dan dirujuk seperti
\cref{lst:g-python}.

\begin{lstlisting}[language=Python, caption={Contoh kode Python}, label=lst:g-python]
def hitung_rerata(nilai):
    """Menghitung rata-rata dari sebuah senarai angka."""
    if not nilai:
        return 0
    return sum(nilai) / len(nilai)
\end{lstlisting}

\begin{lstlisting}[language=SQL, caption={Contoh kueri SQL}, label=lst:g-sql]
SELECT d.nama       AS dimensi,
       AVG(j.kinerja)     AS rerata_kinerja,
       AVG(j.kepentingan) AS rerata_penting
FROM   jawaban j
JOIN   butir b   ON b.id = j.butir_id
JOIN   dimensi d ON d.id = b.dimensi_id
GROUP  BY d.nama
ORDER  BY rerata_kinerja ASC;
\end{lstlisting}

\begin{lstlisting}[language=Java, caption={Contoh kode Java}, label=lst:g-java]
public class Responden {
    private final String kode;
    private final int[] jawaban;

    public Responden(String kode, int[] jawaban) {
        this.kode = kode;
        this.jawaban = jawaban;
    }

    public double rerata() {
        return Arrays.stream(jawaban)
                     .average()
                     .orElse(0.0);
    }
}
\end{lstlisting}

\begin{lstlisting}[language=PHP, caption={Contoh kode PHP}, label=lst:g-php]
<?php
function hitungGap(array $kin, array $pen): array {
    $gap = [];
    foreach ($kin as $b => $nilai) {
        $gap[$b] = round($nilai - $pen[$b], 2);
    }
    return $gap;
}
\end{lstlisting}

\subsection{Kode dibaca dari berkas}
Cara ini lebih dianjurkan untuk lampiran kode karena isinya selalu sinkron
dengan berkas sumber, dapat dicari, dan tidak membengkakkan ukuran berkas
seperti halnya tangkapan layar. Lihat \cref{lst:g-berkas}.

\lstinputlisting[language=Python, caption={Kode yang dibaca langsung dari berkas \texttt{kode/hitung-sus.py}}, label=lst:g-berkas]{kode/hitung-sus.py}

Sebagian baris saja juga dapat ditampilkan dengan parameter
\texttt{firstline} dan \texttt{lastline}.

\lstinputlisting[language=Python, firstline=7, lastline=18, caption={Cuplikan baris 7--18 dari berkas yang sama}, label=lst:g-cuplikan]{kode/hitung-sus.py}

\subsection{Pseudocode}
Untuk menuliskan algoritma, gunakan \texttt{algorithm2e}. Bentuk ini lebih
disukai penguji daripada kode nyata karena tidak terikat bahasa pemrograman
tertentu. Lihat \cref{alg:g-ipa}.

\begin{algorithm}[H]
\caption{Pemetaan kuadran \textit{Importance-Performance Analysis}}
\label{alg:g-ipa}
\KwMasukan{Senarai butir $B$, skor kepentingan $I$, skor kinerja $P$}
\KwKeluaran{Kuadran setiap butir}
$\bar{I} \leftarrow$ rata-rata seluruh nilai $I$\;
$\bar{P} \leftarrow$ rata-rata seluruh nilai $P$\;
\ForEach{butir $b \in B$}{
    \uIf{$I_b \geq \bar{I}$ \textbf{dan} $P_b < \bar{P}$}{
        $K_b \leftarrow$ ``Prioritas utama''\;
    }
    \uElseIf{$I_b \geq \bar{I}$}{
        $K_b \leftarrow$ ``Pertahankan prestasi''\;
    }
    \uElseIf{$P_b < \bar{P}$}{
        $K_b \leftarrow$ ``Prioritas rendah''\;
    }
    \Else{
        $K_b \leftarrow$ ``Berlebihan''\;
    }
}
\Return $K$\;
\end{algorithm}

% ==================================================================
\section{Sitasi dan Daftar Pustaka}
% ==================================================================

Template ini memakai gaya APA nama-tahun. Gaya diatur melalui perintah
\texttt{\textbackslash gayasitasi} pada berkas \texttt{a0-identitas.tex}.

\subsection{Bentuk sitasi}

\begin{daftarpoin}
    \item \texttt{\textbackslash citep\{kunci\}} menghasilkan sitasi dalam
          kurung: \citep{wijaya2022}. Dipakai bila rujukan bukan bagian dari
          kalimat.
    \item \texttt{\textbackslash citet\{kunci\}} menghasilkan sitasi naratif:
          \citet{wijaya2022}. Dipakai bila nama penulis menjadi subjek kalimat.
    \item Sitasi jamak ditulis dalam satu perintah:
          \citep{wijaya2022, pratama2023, sari2024}.
    \item Sitasi dengan halaman:
          \citep[hlm.~118]{barnes2002}.
    \item Sitasi dengan keterangan awalan:
          \citep[lihat][]{martilla1977}.
    \item Menyebut tahun saja: \citeyearpar{martilla1977}.
    \item Menyebut nama penulis saja: \citeauthor{barnes2002}.
\end{daftarpoin}

\subsection{Kutipan langsung}
Kutipan pendek ditulis menyatu dengan paragraf dan diapit tanda kutip.
Kutipan langsung yang panjang, yaitu lebih dari 40 kata, ditulis sebagai blok
tersendiri yang menjorok tanpa tanda kutip, seperti berikut.

\begin{quote}
    \noindent Tulis kutipan panjang Anda di sini. Blok kutipan ditulis menjorok
    ke dalam, tanpa tanda kutip, dan tetap diakhiri sitasi yang memuat nomor
    halaman sumbernya. Kutipan langsung sebaiknya digunakan seperlunya saja;
    utamakan parafrasa dengan bahasa Anda sendiri agar naskah tidak menjadi
    rangkaian kutipan. \citep[hlm.~120]{barnes2002}
\end{quote}

\subsection{Tipe entri pada berkas .bib}
Berkas \texttt{a7-pustaka.bib} memuat katalog seluruh tipe entri yang lazim
dipakai, yaitu \texttt{@article} untuk artikel jurnal, \texttt{@book} untuk
buku, \texttt{@incollection} untuk bab dalam bunga rampai,
\texttt{@inproceedings} untuk prosiding, \texttt{@mastersthesis} untuk skripsi
atau tesis, \texttt{@techreport} untuk laporan lembaga, \texttt{@manual} untuk
standar dan dokumentasi, serta \texttt{@misc} untuk sumber daring. Salin tipe
yang dibutuhkan lalu isi dengan data sebenarnya.

% Perintah \nocite berikut memaksa seluruh tipe entri ikut tampil pada
% Daftar Pustaka meskipun tidak disitasi di dalam teks, sehingga Anda
% dapat melihat bentuk tampilannya. Entri ini akan hilang sendiri
% begitu galeri dimatikan.
\nocite{contoh-buku, contoh-babbuku, contoh-prosiding, contoh-skripsi,
        contoh-laporan, contoh-standar, contoh-web}

Bentuk tampilan masing-masing tipe dapat Anda lihat langsung pada Daftar
Pustaka naskah ini.

% ==================================================================
\section{Rujukan Silang}
% ==================================================================

Seluruh rujukan silang sebaiknya memakai \texttt{\textbackslash cref} agar kata
``Gambar'', ``Tabel'', atau ``Persamaan'' muncul otomatis dan penomorannya
selalu sinkron: \cref{fig:g-tunggal}, \cref{tab:g-biasa}, \cref{eq:g-rerata},
\cref{lst:g-python}, dan \cref{alg:g-ipa}.

Beberapa objek sekaligus dapat dirujuk dalam satu perintah dan LaTeX akan
menyusun kata sambungnya sendiri: \cref{tab:g-biasa,tab:g-booktabs,tab:g-gabung}.

Rujukan ke rumusan masalah juga tersedia. Perintah \texttt{\textbackslash
cref\{rq:kualitas\}} menghasilkan \cref{rq:kualitas}, sehingga penomoran RQ di
seluruh naskah menyesuaikan sendiri bila urutannya diubah.
 di laporan.tex
% sebelum naskah dikumpulkan. Berkas ini tidak akan ikut tercetak,
% tetapi tetap tersimpan sebagai rujukan.
%==================================================================

% Galeri disertakan setelah lampiran, sehingga tanpa pengaturan ini
% penomorannya akan ikut melanjutkan huruf lampiran (D, E, dst.).
% Blok berikut membuatnya bernomor "G" agar jelas terpisah.
\clearpage
\setcounter{chapter}{0}
\setcounter{figure}{0}
\setcounter{table}{0}
\setcounter{equation}{0}
\renewcommand{\thechapter}{G}
\renewcommand{\thesection}{G.\arabic{section}}
\renewcommand{\thesubsection}{\thesection.\arabic{subsection}}
\renewcommand{\thefigure}{G.\arabic{figure}}
\renewcommand{\thetable}{G.\arabic{table}}
\renewcommand{\theequation}{G.\arabic{equation}}

\chapter*{GALERI ELEMEN (HAPUS SEBELUM DIKUMPULKAN)}
\addcontentsline{toc}{chapter}{GALERI ELEMEN (HAPUS SEBELUM DIKUMPULKAN)}

Bab ini bukan bagian dari naskah Tugas Akhir. Isinya adalah contoh siap salin
untuk seluruh elemen yang lazim dibutuhkan. Matikan dengan memberi tanda
\texttt{\%} pada baris \texttt{\textbackslash include\{zz-galeri-elemen\}} di
berkas \texttt{laporan.tex}.

% ==================================================================
\section{Tabel}
% ==================================================================

\subsection{Tabel biasa}
Keterangan tabel diletakkan di \textbf{atas} tabel. Setiap tabel wajib diberi
\texttt{\textbackslash label} dan disebut dalam teks, misalnya
\cref{tab:g-biasa}.

\begin{table}[H]
    \centering
    \caption{Contoh tabel biasa}
    \label{tab:g-biasa}
    \begin{tabular}{|c|l|c|}
        \hline
        \textbf{No} & \textbf{Keterangan} & \textbf{Nilai} \\ \hline
        1 & Baris pertama & 3,12 \\ \hline
        2 & Baris kedua   & 3,68 \\ \hline
        3 & Baris ketiga  & 3,54 \\ \hline
    \end{tabular}
\end{table}

\subsection{Tabel dengan garis gaya ilmiah (booktabs)}
Sebagian program studi menghendaki tabel tanpa garis vertikal.

\begin{table}[H]
    \centering
    \caption{Contoh tabel gaya booktabs}
    \label{tab:g-booktabs}
    \begin{tabular}{lcc}
        \toprule
        \textbf{Dimensi} & \textbf{Kepentingan} & \textbf{Kinerja} \\
        \midrule
        \textit{Usability}           & 4,41 & 3,12 \\
        \textit{Information quality} & 4,22 & 3,68 \\
        \bottomrule
    \end{tabular}
\end{table}

\subsection{Tabel bersambung antarhalaman (longtable)}
Gunakan \texttt{longtable} bila tabel lebih panjang daripada satu halaman.
Baris kepala akan diulang otomatis, dan halaman sambungan diberi keterangan
``Lanjutan'' sesuai kelaziman naskah Indonesia.

\begin{longtable}{|c|p{6.5cm}|c|c|}
    \caption{Contoh tabel bersambung antarhalaman}
    \label{tab:g-longtable} \\
    \hline
    \textbf{No} & \textbf{Pernyataan} & \textbf{Kepentingan} & \textbf{Kinerja} \\ \hline
    \endfirsthead

    \multicolumn{4}{l}{\footnotesize\textit{Lanjutan \cref{tab:g-longtable}}} \\
    \hline
    \textbf{No} & \textbf{Pernyataan} & \textbf{Kepentingan} & \textbf{Kinerja} \\ \hline
    \endhead

    \hline
    \multicolumn{4}{r}{\footnotesize\textit{Bersambung ke halaman berikutnya}} \\
    \endfoot

    \hline
    \endlastfoot

    1  & Saya merasa sistem mudah dipelajari & 4,44 & 3,10 \\ \hline
    2  & Interaksi dengan sistem jelas dan mudah dipahami & 4,38 & 3,21 \\ \hline
    3  & Sistem mudah dioperasikan & 4,41 & 3,15 \\ \hline
    4  & Navigasi antarhalaman terasa mudah & 4,61 & 3,03 \\ \hline
    5  & Tampilan sistem sesuai dengan jenis layanannya & 4,02 & 3,44 \\ \hline
    6  & Tampilan sistem menarik & 3,88 & 3,52 \\ \hline
    7  & Sistem menyampaikan kesan kompeten & 4,12 & 3,40 \\ \hline
    8  & Sistem menciptakan pengalaman yang menyenangkan & 3,95 & 3,33 \\ \hline
    9  & Informasi yang disajikan akurat & 4,30 & 3,71 \\ \hline
    10 & Informasi yang disajikan tepercaya & 4,28 & 3,69 \\ \hline
    11 & Informasi selalu diperbarui & 4,19 & 3,58 \\ \hline
    12 & Informasi relevan dengan kebutuhan & 4,25 & 3,66 \\ \hline
    13 & Informasi mudah dipahami & 4,21 & 3,74 \\ \hline
    14 & Informasi disajikan pada tingkat rincian yang tepat & 4,08 & 3,70 \\ \hline
    15 & Informasi disajikan dalam format yang sesuai & 4,15 & 3,68 \\ \hline
    16 & Sistem memberikan rasa aman & 4,09 & 3,61 \\ \hline
    17 & Data pribadi terasa terlindungi & 4,18 & 3,55 \\ \hline
    18 & Sistem memudahkan komunikasi dengan pengelola & 3,97 & 3,42 \\ \hline
    19 & Sistem menciptakan rasa keterhubungan & 3,82 & 3,39 \\ \hline
    20 & Layanan diberikan sesuai yang dijanjikan & 4,11 & 3,58 \\ \hline
\end{longtable}

\subsection{Tabel dengan sel gabung (multirow dan multicolumn)}

\begin{table}[H]
    \centering
    \caption{Contoh tabel dengan sel gabung}
    \label{tab:g-gabung}
    \begin{tabular}{|c|c|c|c|c|}
        \hline
        \multirow{2}{*}{\textbf{Butir}} & \multicolumn{2}{c|}{\textbf{Sebelum}} & \multicolumn{2}{c|}{\textbf{Sesudah}} \\ \cline{2-5}
         & \textbf{Rerata} & \textbf{SB} & \textbf{Rerata} & \textbf{SB} \\ \hline
        B1 & 3,10 & 0,82 & 4,05 & 0,61 \\ \hline
        B2 & 3,22 & 0,77 & 4,12 & 0,58 \\ \hline
    \end{tabular}
\end{table}

\subsection{Tabel berkolom otomatis (tabularx)}
Lebar kolom bertipe \texttt{X} menyesuaikan sendiri agar tabel tepat memenuhi
lebar teks.

\begin{table}[H]
    \centering
    \caption{Contoh tabel dengan lebar kolom otomatis}
    \label{tab:g-tabularx}
    \begin{tabularx}{\textwidth}{|c|X|X|}
        \hline
        \textbf{No} & \textbf{Temuan} & \textbf{Tindak lanjut} \\ \hline
        1 & Menu utama tidak konsisten antarhalaman sehingga pengguna kehilangan orientasi. & Menyeragamkan struktur menu pada seluruh halaman. \\ \hline
        2 & Pesan kesalahan tidak menjelaskan penyebab masalah. & Menulis ulang pesan kesalahan dengan bahasa yang spesifik. \\ \hline
    \end{tabularx}
\end{table}

\subsection{Tabel dengan catatan kaki (threeparttable)}
Gunakan bila tabel memerlukan keterangan simbol atau sumber.

\begin{table}[H]
    \centering
    \begin{threeparttable}
        \caption{Contoh tabel dengan catatan kaki}
        \label{tab:g-catatan}
        \begin{tabular}{|l|c|c|}
            \hline
            \textbf{Dimensi} & \textbf{Rerata} & \textbf{Kategori} \\ \hline
            \textit{Usability}\tnote{a}           & 3,12 & Cukup \\ \hline
            \textit{Information quality}          & 3,68 & Baik  \\ \hline
            \textit{Service interaction}\tnote{b} & 3,54 & Baik  \\ \hline
        \end{tabular}
        \begin{tablenotes}\footnotesize
            \item[a] Terdiri atas 8 butir pernyataan.
            \item[b] Terdiri atas 7 butir pernyataan.
            \item Sumber: Hasil pengolahan data penelitian (2026).
        \end{tablenotes}
    \end{threeparttable}
\end{table}

\subsection{Tabel yang terlalu lebar}
Bila tabel sedikit melebihi lebar teks, perkecil jarak antarkolom dengan
\texttt{\textbackslash tabcolsep} atau bungkus dengan \texttt{adjustbox}.

\begin{table}[H]
    \centering
    \caption{Contoh tabel lebar yang dikecilkan otomatis}
    \label{tab:g-lebar}
    \footnotesize
    \begin{adjustbox}{max width=\textwidth}
    \begin{tabular}{|c|c|c|c|c|c|c|c|c|c|c|}
        \hline
        \textbf{Butir} & \textbf{R1} & \textbf{R2} & \textbf{R3} & \textbf{R4} & \textbf{R5} & \textbf{R6} & \textbf{R7} & \textbf{R8} & \textbf{Rerata} & \textbf{Simpangan Baku} \\ \hline
        B1 & 4 & 5 & 4 & 3 & 4 & 5 & 4 & 4 & 4,13 & 0,64 \\ \hline
        B2 & 3 & 4 & 4 & 4 & 5 & 4 & 3 & 4 & 3,88 & 0,64 \\ \hline
    \end{tabular}
    \end{adjustbox}
\end{table}

\subsection{Tabel lanskap}
Paket \texttt{pdflscape} ikut memutar halaman pada berkas PDF, sehingga pembaca
tidak perlu memiringkan kepala.

\begin{landscape}
\begin{table}[H]
    \centering
    \caption{Contoh tabel pada halaman lanskap}
    \label{tab:g-lanskap}
    \footnotesize
    \begin{tabular}{|c|p{4cm}|p{4cm}|p{4cm}|p{4cm}|}
        \hline
        \textbf{No} & \textbf{Kondisi berjalan} & \textbf{Masalah} & \textbf{Usulan} & \textbf{Dampak yang diharapkan} \\ \hline
        1 & Menu tersebar pada tiga tingkat & Pengguna kesulitan menemukan menu & Menyusun ulang menjadi dua tingkat & Waktu pencarian menu berkurang \\ \hline
        2 & Alur pengisian tidak menampilkan tahapan & Pengguna tidak tahu posisinya & Menambahkan penanda tahapan & Kesalahan pengisian berkurang \\ \hline
    \end{tabular}
\end{table}
\end{landscape}

% ==================================================================
\section{Gambar}
% ==================================================================

\subsection{Gambar tunggal}
Keterangan gambar diletakkan di \textbf{bawah} gambar. Lihat
\cref{fig:g-tunggal}.

\begin{figure}[H]
    \centering
    \includegraphics[width=0.6\linewidth]{contoh-diagram}
    \caption{Contoh gambar tunggal}
    \label{fig:g-tunggal}
\end{figure}

\subsection{Gambar dengan keterangan sumber}
Gunakan perintah \texttt{\textbackslash sumber} bila gambar dikutip dari
sumber lain.

\begin{figure}[H]
    \centering
    \includegraphics[width=0.55\linewidth]{contoh-diagram}
    \caption{Contoh gambar yang dikutip dari sumber lain}
    \label{fig:g-sumber}
\end{figure}
\sumber{\citet{barnes2002}}

\subsection{Dua gambar berdampingan}
Gunakan \texttt{subcaption} agar tiap bagian memperoleh label (a) dan (b) yang
dapat dirujuk terpisah, misalnya \cref{fig:g-kiri} dan \cref{fig:g-kanan}.

\begin{figure}[H]
    \centering
    \begin{subfigure}[t]{0.47\linewidth}
        \centering
        \includegraphics[width=\linewidth]{contoh-tangkapan-a}
        \caption{Tampilan sebelum perbaikan}
        \label{fig:g-kiri}
    \end{subfigure}
    \hfill
    \begin{subfigure}[t]{0.47\linewidth}
        \centering
        \includegraphics[width=\linewidth]{contoh-tangkapan-b}
        \caption{Tampilan sesudah perbaikan}
        \label{fig:g-kanan}
    \end{subfigure}
    \caption{Perbandingan tampilan antarmuka}
    \label{fig:g-berdampingan}
\end{figure}

\subsection{Gambar tinggi satu halaman}
Untuk gambar yang tinggi, seperti diagram alir panjang, gunakan
\texttt{height} agar tidak meluber ke halaman berikutnya.

\begin{figure}[H]
    \centering
    \includegraphics[height=0.78\textheight,keepaspectratio]{contoh-tinggi}
    \caption{Contoh gambar tinggi satu halaman}
    \label{fig:g-tinggi}
\end{figure}

\subsection{Gambar lanskap}

\begin{landscape}
\begin{figure}[H]
    \centering
    \includegraphics[width=0.92\linewidth]{contoh-lanskap}
    \caption{Contoh gambar pada halaman lanskap}
    \label{fig:g-lanskap}
\end{figure}
\end{landscape}

\subsection{Catatan penempatan gambar}
Parameter \texttt{[H]} memaksa gambar muncul tepat di posisi penulisannya.
Parameter \texttt{[htbp]} memberi LaTeX kebebasan menempatkan gambar pada
posisi terbaik, sehingga halaman tidak menyisakan ruang kosong. Gunakan
\texttt{[H]} bila urutan gambar penting, dan \texttt{[htbp]} bila
mengutamakan kerapian halaman.

% ==================================================================
\section{Persamaan}
% ==================================================================

Persamaan dalam baris teks ditulis seperti $\bar{x} = \sum x_i / n$, sedangkan
persamaan bernomor ditulis terpisah seperti \cref{eq:g-rerata}.

\begin{equation}
    \label{eq:g-rerata}
    \bar{x} = \frac{1}{n}\sum_{i=1}^{n} x_{i}
\end{equation}

Persamaan bertingkat ditulis dengan lingkungan \texttt{align}, dan setiap baris
memperoleh nomor tersendiri seperti \cref{eq:g-gap} dan \cref{eq:g-tki}.

\begin{align}
    \text{Gap}_i &= \bar{P}_i - \bar{I}_i \label{eq:g-gap} \\
    \text{TKI}_i &= \frac{\bar{P}_i}{\bar{I}_i} \times 100\% \label{eq:g-tki}
\end{align}

Beberapa bentuk yang sering dibutuhkan: pecahan $\frac{a}{b}$, akar
$\sqrt[3]{x}$, jumlah $\sum_{i=1}^{n}$, hasil kali $\prod_{i=1}^{n}$, integral
$\int_{0}^{1}$, huruf Yunani $\alpha, \beta, \sigma, \mu, \kappa$, serta matriks
berikut.

\begin{equation}
    \label{eq:g-matriks}
    \mathbf{D} =
    \begin{bmatrix}
        d_{11} & d_{12} & \cdots & d_{1n} \\
        d_{21} & d_{22} & \cdots & d_{2n} \\
        \vdots & \vdots & \ddots & \vdots \\
        d_{m1} & d_{m2} & \cdots & d_{mn}
    \end{bmatrix}
\end{equation}

% ==================================================================
\section{Kode Program}
% ==================================================================

\subsection{Kode ditulis langsung}
Kode dapat ditulis langsung di dalam naskah dan dirujuk seperti
\cref{lst:g-python}.

\begin{lstlisting}[language=Python, caption={Contoh kode Python}, label=lst:g-python]
def hitung_rerata(nilai):
    """Menghitung rata-rata dari sebuah senarai angka."""
    if not nilai:
        return 0
    return sum(nilai) / len(nilai)
\end{lstlisting}

\begin{lstlisting}[language=SQL, caption={Contoh kueri SQL}, label=lst:g-sql]
SELECT d.nama       AS dimensi,
       AVG(j.kinerja)     AS rerata_kinerja,
       AVG(j.kepentingan) AS rerata_penting
FROM   jawaban j
JOIN   butir b   ON b.id = j.butir_id
JOIN   dimensi d ON d.id = b.dimensi_id
GROUP  BY d.nama
ORDER  BY rerata_kinerja ASC;
\end{lstlisting}

\begin{lstlisting}[language=Java, caption={Contoh kode Java}, label=lst:g-java]
public class Responden {
    private final String kode;
    private final int[] jawaban;

    public Responden(String kode, int[] jawaban) {
        this.kode = kode;
        this.jawaban = jawaban;
    }

    public double rerata() {
        return Arrays.stream(jawaban)
                     .average()
                     .orElse(0.0);
    }
}
\end{lstlisting}

\begin{lstlisting}[language=PHP, caption={Contoh kode PHP}, label=lst:g-php]
<?php
function hitungGap(array $kin, array $pen): array {
    $gap = [];
    foreach ($kin as $b => $nilai) {
        $gap[$b] = round($nilai - $pen[$b], 2);
    }
    return $gap;
}
\end{lstlisting}

\subsection{Kode dibaca dari berkas}
Cara ini lebih dianjurkan untuk lampiran kode karena isinya selalu sinkron
dengan berkas sumber, dapat dicari, dan tidak membengkakkan ukuran berkas
seperti halnya tangkapan layar. Lihat \cref{lst:g-berkas}.

\lstinputlisting[language=Python, caption={Kode yang dibaca langsung dari berkas \texttt{kode/hitung-sus.py}}, label=lst:g-berkas]{kode/hitung-sus.py}

Sebagian baris saja juga dapat ditampilkan dengan parameter
\texttt{firstline} dan \texttt{lastline}.

\lstinputlisting[language=Python, firstline=7, lastline=18, caption={Cuplikan baris 7--18 dari berkas yang sama}, label=lst:g-cuplikan]{kode/hitung-sus.py}

\subsection{Pseudocode}
Untuk menuliskan algoritma, gunakan \texttt{algorithm2e}. Bentuk ini lebih
disukai penguji daripada kode nyata karena tidak terikat bahasa pemrograman
tertentu. Lihat \cref{alg:g-ipa}.

\begin{algorithm}[H]
\caption{Pemetaan kuadran \textit{Importance-Performance Analysis}}
\label{alg:g-ipa}
\KwMasukan{Senarai butir $B$, skor kepentingan $I$, skor kinerja $P$}
\KwKeluaran{Kuadran setiap butir}
$\bar{I} \leftarrow$ rata-rata seluruh nilai $I$\;
$\bar{P} \leftarrow$ rata-rata seluruh nilai $P$\;
\ForEach{butir $b \in B$}{
    \uIf{$I_b \geq \bar{I}$ \textbf{dan} $P_b < \bar{P}$}{
        $K_b \leftarrow$ ``Prioritas utama''\;
    }
    \uElseIf{$I_b \geq \bar{I}$}{
        $K_b \leftarrow$ ``Pertahankan prestasi''\;
    }
    \uElseIf{$P_b < \bar{P}$}{
        $K_b \leftarrow$ ``Prioritas rendah''\;
    }
    \Else{
        $K_b \leftarrow$ ``Berlebihan''\;
    }
}
\Return $K$\;
\end{algorithm}

% ==================================================================
\section{Sitasi dan Daftar Pustaka}
% ==================================================================

Template ini memakai gaya APA nama-tahun. Gaya diatur melalui perintah
\texttt{\textbackslash gayasitasi} pada berkas \texttt{a0-identitas.tex}.

\subsection{Bentuk sitasi}

\begin{daftarpoin}
    \item \texttt{\textbackslash citep\{kunci\}} menghasilkan sitasi dalam
          kurung: \citep{wijaya2022}. Dipakai bila rujukan bukan bagian dari
          kalimat.
    \item \texttt{\textbackslash citet\{kunci\}} menghasilkan sitasi naratif:
          \citet{wijaya2022}. Dipakai bila nama penulis menjadi subjek kalimat.
    \item Sitasi jamak ditulis dalam satu perintah:
          \citep{wijaya2022, pratama2023, sari2024}.
    \item Sitasi dengan halaman:
          \citep[hlm.~118]{barnes2002}.
    \item Sitasi dengan keterangan awalan:
          \citep[lihat][]{martilla1977}.
    \item Menyebut tahun saja: \citeyearpar{martilla1977}.
    \item Menyebut nama penulis saja: \citeauthor{barnes2002}.
\end{daftarpoin}

\subsection{Kutipan langsung}
Kutipan pendek ditulis menyatu dengan paragraf dan diapit tanda kutip.
Kutipan langsung yang panjang, yaitu lebih dari 40 kata, ditulis sebagai blok
tersendiri yang menjorok tanpa tanda kutip, seperti berikut.

\begin{quote}
    \noindent Tulis kutipan panjang Anda di sini. Blok kutipan ditulis menjorok
    ke dalam, tanpa tanda kutip, dan tetap diakhiri sitasi yang memuat nomor
    halaman sumbernya. Kutipan langsung sebaiknya digunakan seperlunya saja;
    utamakan parafrasa dengan bahasa Anda sendiri agar naskah tidak menjadi
    rangkaian kutipan. \citep[hlm.~120]{barnes2002}
\end{quote}

\subsection{Tipe entri pada berkas .bib}
Berkas \texttt{a7-pustaka.bib} memuat katalog seluruh tipe entri yang lazim
dipakai, yaitu \texttt{@article} untuk artikel jurnal, \texttt{@book} untuk
buku, \texttt{@incollection} untuk bab dalam bunga rampai,
\texttt{@inproceedings} untuk prosiding, \texttt{@mastersthesis} untuk skripsi
atau tesis, \texttt{@techreport} untuk laporan lembaga, \texttt{@manual} untuk
standar dan dokumentasi, serta \texttt{@misc} untuk sumber daring. Salin tipe
yang dibutuhkan lalu isi dengan data sebenarnya.

% Perintah \nocite berikut memaksa seluruh tipe entri ikut tampil pada
% Daftar Pustaka meskipun tidak disitasi di dalam teks, sehingga Anda
% dapat melihat bentuk tampilannya. Entri ini akan hilang sendiri
% begitu galeri dimatikan.
\nocite{contoh-buku, contoh-babbuku, contoh-prosiding, contoh-skripsi,
        contoh-laporan, contoh-standar, contoh-web}

Bentuk tampilan masing-masing tipe dapat Anda lihat langsung pada Daftar
Pustaka naskah ini.

% ==================================================================
\section{Rujukan Silang}
% ==================================================================

Seluruh rujukan silang sebaiknya memakai \texttt{\textbackslash cref} agar kata
``Gambar'', ``Tabel'', atau ``Persamaan'' muncul otomatis dan penomorannya
selalu sinkron: \cref{fig:g-tunggal}, \cref{tab:g-biasa}, \cref{eq:g-rerata},
\cref{lst:g-python}, dan \cref{alg:g-ipa}.

Beberapa objek sekaligus dapat dirujuk dalam satu perintah dan LaTeX akan
menyusun kata sambungnya sendiri: \cref{tab:g-biasa,tab:g-booktabs,tab:g-gabung}.

Rujukan ke rumusan masalah juga tersedia. Perintah \texttt{\textbackslash
cref\{rq:kualitas\}} menghasilkan \cref{rq:kualitas}, sehingga penomoran RQ di
seluruh naskah menyesuaikan sendiri bila urutannya diubah.
 di laporan.tex
% sebelum naskah dikumpulkan. Berkas ini tidak akan ikut tercetak,
% tetapi tetap tersimpan sebagai rujukan.
%==================================================================

% Galeri disertakan setelah lampiran, sehingga tanpa pengaturan ini
% penomorannya akan ikut melanjutkan huruf lampiran (D, E, dst.).
% Blok berikut membuatnya bernomor "G" agar jelas terpisah.
\clearpage
\setcounter{chapter}{0}
\setcounter{figure}{0}
\setcounter{table}{0}
\setcounter{equation}{0}
\renewcommand{\thechapter}{G}
\renewcommand{\thesection}{G.\arabic{section}}
\renewcommand{\thesubsection}{\thesection.\arabic{subsection}}
\renewcommand{\thefigure}{G.\arabic{figure}}
\renewcommand{\thetable}{G.\arabic{table}}
\renewcommand{\theequation}{G.\arabic{equation}}

\chapter*{GALERI ELEMEN (HAPUS SEBELUM DIKUMPULKAN)}
\addcontentsline{toc}{chapter}{GALERI ELEMEN (HAPUS SEBELUM DIKUMPULKAN)}

Bab ini bukan bagian dari naskah Tugas Akhir. Isinya adalah contoh siap salin
untuk seluruh elemen yang lazim dibutuhkan. Matikan dengan memberi tanda
\texttt{\%} pada baris \texttt{\textbackslash include\{zz-galeri-elemen\}} di
berkas \texttt{laporan.tex}.

% ==================================================================
\section{Tabel}
% ==================================================================

\subsection{Tabel biasa}
Keterangan tabel diletakkan di \textbf{atas} tabel. Setiap tabel wajib diberi
\texttt{\textbackslash label} dan disebut dalam teks, misalnya
\cref{tab:g-biasa}.

\begin{table}[H]
    \centering
    \caption{Contoh tabel biasa}
    \label{tab:g-biasa}
    \begin{tabular}{|c|l|c|}
        \hline
        \textbf{No} & \textbf{Keterangan} & \textbf{Nilai} \\ \hline
        1 & Baris pertama & 3,12 \\ \hline
        2 & Baris kedua   & 3,68 \\ \hline
        3 & Baris ketiga  & 3,54 \\ \hline
    \end{tabular}
\end{table}

\subsection{Tabel dengan garis gaya ilmiah (booktabs)}
Sebagian program studi menghendaki tabel tanpa garis vertikal.

\begin{table}[H]
    \centering
    \caption{Contoh tabel gaya booktabs}
    \label{tab:g-booktabs}
    \begin{tabular}{lcc}
        \toprule
        \textbf{Dimensi} & \textbf{Kepentingan} & \textbf{Kinerja} \\
        \midrule
        \textit{Usability}           & 4,41 & 3,12 \\
        \textit{Information quality} & 4,22 & 3,68 \\
        \bottomrule
    \end{tabular}
\end{table}

\subsection{Tabel bersambung antarhalaman (longtable)}
Gunakan \texttt{longtable} bila tabel lebih panjang daripada satu halaman.
Baris kepala akan diulang otomatis, dan halaman sambungan diberi keterangan
``Lanjutan'' sesuai kelaziman naskah Indonesia.

\begin{longtable}{|c|p{6.5cm}|c|c|}
    \caption{Contoh tabel bersambung antarhalaman}
    \label{tab:g-longtable} \\
    \hline
    \textbf{No} & \textbf{Pernyataan} & \textbf{Kepentingan} & \textbf{Kinerja} \\ \hline
    \endfirsthead

    \multicolumn{4}{l}{\footnotesize\textit{Lanjutan \cref{tab:g-longtable}}} \\
    \hline
    \textbf{No} & \textbf{Pernyataan} & \textbf{Kepentingan} & \textbf{Kinerja} \\ \hline
    \endhead

    \hline
    \multicolumn{4}{r}{\footnotesize\textit{Bersambung ke halaman berikutnya}} \\
    \endfoot

    \hline
    \endlastfoot

    1  & Saya merasa sistem mudah dipelajari & 4,44 & 3,10 \\ \hline
    2  & Interaksi dengan sistem jelas dan mudah dipahami & 4,38 & 3,21 \\ \hline
    3  & Sistem mudah dioperasikan & 4,41 & 3,15 \\ \hline
    4  & Navigasi antarhalaman terasa mudah & 4,61 & 3,03 \\ \hline
    5  & Tampilan sistem sesuai dengan jenis layanannya & 4,02 & 3,44 \\ \hline
    6  & Tampilan sistem menarik & 3,88 & 3,52 \\ \hline
    7  & Sistem menyampaikan kesan kompeten & 4,12 & 3,40 \\ \hline
    8  & Sistem menciptakan pengalaman yang menyenangkan & 3,95 & 3,33 \\ \hline
    9  & Informasi yang disajikan akurat & 4,30 & 3,71 \\ \hline
    10 & Informasi yang disajikan tepercaya & 4,28 & 3,69 \\ \hline
    11 & Informasi selalu diperbarui & 4,19 & 3,58 \\ \hline
    12 & Informasi relevan dengan kebutuhan & 4,25 & 3,66 \\ \hline
    13 & Informasi mudah dipahami & 4,21 & 3,74 \\ \hline
    14 & Informasi disajikan pada tingkat rincian yang tepat & 4,08 & 3,70 \\ \hline
    15 & Informasi disajikan dalam format yang sesuai & 4,15 & 3,68 \\ \hline
    16 & Sistem memberikan rasa aman & 4,09 & 3,61 \\ \hline
    17 & Data pribadi terasa terlindungi & 4,18 & 3,55 \\ \hline
    18 & Sistem memudahkan komunikasi dengan pengelola & 3,97 & 3,42 \\ \hline
    19 & Sistem menciptakan rasa keterhubungan & 3,82 & 3,39 \\ \hline
    20 & Layanan diberikan sesuai yang dijanjikan & 4,11 & 3,58 \\ \hline
\end{longtable}

\subsection{Tabel dengan sel gabung (multirow dan multicolumn)}

\begin{table}[H]
    \centering
    \caption{Contoh tabel dengan sel gabung}
    \label{tab:g-gabung}
    \begin{tabular}{|c|c|c|c|c|}
        \hline
        \multirow{2}{*}{\textbf{Butir}} & \multicolumn{2}{c|}{\textbf{Sebelum}} & \multicolumn{2}{c|}{\textbf{Sesudah}} \\ \cline{2-5}
         & \textbf{Rerata} & \textbf{SB} & \textbf{Rerata} & \textbf{SB} \\ \hline
        B1 & 3,10 & 0,82 & 4,05 & 0,61 \\ \hline
        B2 & 3,22 & 0,77 & 4,12 & 0,58 \\ \hline
    \end{tabular}
\end{table}

\subsection{Tabel berkolom otomatis (tabularx)}
Lebar kolom bertipe \texttt{X} menyesuaikan sendiri agar tabel tepat memenuhi
lebar teks.

\begin{table}[H]
    \centering
    \caption{Contoh tabel dengan lebar kolom otomatis}
    \label{tab:g-tabularx}
    \begin{tabularx}{\textwidth}{|c|X|X|}
        \hline
        \textbf{No} & \textbf{Temuan} & \textbf{Tindak lanjut} \\ \hline
        1 & Menu utama tidak konsisten antarhalaman sehingga pengguna kehilangan orientasi. & Menyeragamkan struktur menu pada seluruh halaman. \\ \hline
        2 & Pesan kesalahan tidak menjelaskan penyebab masalah. & Menulis ulang pesan kesalahan dengan bahasa yang spesifik. \\ \hline
    \end{tabularx}
\end{table}

\subsection{Tabel dengan catatan kaki (threeparttable)}
Gunakan bila tabel memerlukan keterangan simbol atau sumber.

\begin{table}[H]
    \centering
    \begin{threeparttable}
        \caption{Contoh tabel dengan catatan kaki}
        \label{tab:g-catatan}
        \begin{tabular}{|l|c|c|}
            \hline
            \textbf{Dimensi} & \textbf{Rerata} & \textbf{Kategori} \\ \hline
            \textit{Usability}\tnote{a}           & 3,12 & Cukup \\ \hline
            \textit{Information quality}          & 3,68 & Baik  \\ \hline
            \textit{Service interaction}\tnote{b} & 3,54 & Baik  \\ \hline
        \end{tabular}
        \begin{tablenotes}\footnotesize
            \item[a] Terdiri atas 8 butir pernyataan.
            \item[b] Terdiri atas 7 butir pernyataan.
            \item Sumber: Hasil pengolahan data penelitian (2026).
        \end{tablenotes}
    \end{threeparttable}
\end{table}

\subsection{Tabel yang terlalu lebar}
Bila tabel sedikit melebihi lebar teks, perkecil jarak antarkolom dengan
\texttt{\textbackslash tabcolsep} atau bungkus dengan \texttt{adjustbox}.

\begin{table}[H]
    \centering
    \caption{Contoh tabel lebar yang dikecilkan otomatis}
    \label{tab:g-lebar}
    \footnotesize
    \begin{adjustbox}{max width=\textwidth}
    \begin{tabular}{|c|c|c|c|c|c|c|c|c|c|c|}
        \hline
        \textbf{Butir} & \textbf{R1} & \textbf{R2} & \textbf{R3} & \textbf{R4} & \textbf{R5} & \textbf{R6} & \textbf{R7} & \textbf{R8} & \textbf{Rerata} & \textbf{Simpangan Baku} \\ \hline
        B1 & 4 & 5 & 4 & 3 & 4 & 5 & 4 & 4 & 4,13 & 0,64 \\ \hline
        B2 & 3 & 4 & 4 & 4 & 5 & 4 & 3 & 4 & 3,88 & 0,64 \\ \hline
    \end{tabular}
    \end{adjustbox}
\end{table}

\subsection{Tabel lanskap}
Paket \texttt{pdflscape} ikut memutar halaman pada berkas PDF, sehingga pembaca
tidak perlu memiringkan kepala.

\begin{landscape}
\begin{table}[H]
    \centering
    \caption{Contoh tabel pada halaman lanskap}
    \label{tab:g-lanskap}
    \footnotesize
    \begin{tabular}{|c|p{4cm}|p{4cm}|p{4cm}|p{4cm}|}
        \hline
        \textbf{No} & \textbf{Kondisi berjalan} & \textbf{Masalah} & \textbf{Usulan} & \textbf{Dampak yang diharapkan} \\ \hline
        1 & Menu tersebar pada tiga tingkat & Pengguna kesulitan menemukan menu & Menyusun ulang menjadi dua tingkat & Waktu pencarian menu berkurang \\ \hline
        2 & Alur pengisian tidak menampilkan tahapan & Pengguna tidak tahu posisinya & Menambahkan penanda tahapan & Kesalahan pengisian berkurang \\ \hline
    \end{tabular}
\end{table}
\end{landscape}

% ==================================================================
\section{Gambar}
% ==================================================================

\subsection{Gambar tunggal}
Keterangan gambar diletakkan di \textbf{bawah} gambar. Lihat
\cref{fig:g-tunggal}.

\begin{figure}[H]
    \centering
    \includegraphics[width=0.6\linewidth]{contoh-diagram}
    \caption{Contoh gambar tunggal}
    \label{fig:g-tunggal}
\end{figure}

\subsection{Gambar dengan keterangan sumber}
Gunakan perintah \texttt{\textbackslash sumber} bila gambar dikutip dari
sumber lain.

\begin{figure}[H]
    \centering
    \includegraphics[width=0.55\linewidth]{contoh-diagram}
    \caption{Contoh gambar yang dikutip dari sumber lain}
    \label{fig:g-sumber}
\end{figure}
\sumber{\citet{barnes2002}}

\subsection{Dua gambar berdampingan}
Gunakan \texttt{subcaption} agar tiap bagian memperoleh label (a) dan (b) yang
dapat dirujuk terpisah, misalnya \cref{fig:g-kiri} dan \cref{fig:g-kanan}.

\begin{figure}[H]
    \centering
    \begin{subfigure}[t]{0.47\linewidth}
        \centering
        \includegraphics[width=\linewidth]{contoh-tangkapan-a}
        \caption{Tampilan sebelum perbaikan}
        \label{fig:g-kiri}
    \end{subfigure}
    \hfill
    \begin{subfigure}[t]{0.47\linewidth}
        \centering
        \includegraphics[width=\linewidth]{contoh-tangkapan-b}
        \caption{Tampilan sesudah perbaikan}
        \label{fig:g-kanan}
    \end{subfigure}
    \caption{Perbandingan tampilan antarmuka}
    \label{fig:g-berdampingan}
\end{figure}

\subsection{Gambar tinggi satu halaman}
Untuk gambar yang tinggi, seperti diagram alir panjang, gunakan
\texttt{height} agar tidak meluber ke halaman berikutnya.

\begin{figure}[H]
    \centering
    \includegraphics[height=0.78\textheight,keepaspectratio]{contoh-tinggi}
    \caption{Contoh gambar tinggi satu halaman}
    \label{fig:g-tinggi}
\end{figure}

\subsection{Gambar lanskap}

\begin{landscape}
\begin{figure}[H]
    \centering
    \includegraphics[width=0.92\linewidth]{contoh-lanskap}
    \caption{Contoh gambar pada halaman lanskap}
    \label{fig:g-lanskap}
\end{figure}
\end{landscape}

\subsection{Catatan penempatan gambar}
Parameter \texttt{[H]} memaksa gambar muncul tepat di posisi penulisannya.
Parameter \texttt{[htbp]} memberi LaTeX kebebasan menempatkan gambar pada
posisi terbaik, sehingga halaman tidak menyisakan ruang kosong. Gunakan
\texttt{[H]} bila urutan gambar penting, dan \texttt{[htbp]} bila
mengutamakan kerapian halaman.

% ==================================================================
\section{Persamaan}
% ==================================================================

Persamaan dalam baris teks ditulis seperti $\bar{x} = \sum x_i / n$, sedangkan
persamaan bernomor ditulis terpisah seperti \cref{eq:g-rerata}.

\begin{equation}
    \label{eq:g-rerata}
    \bar{x} = \frac{1}{n}\sum_{i=1}^{n} x_{i}
\end{equation}

Persamaan bertingkat ditulis dengan lingkungan \texttt{align}, dan setiap baris
memperoleh nomor tersendiri seperti \cref{eq:g-gap} dan \cref{eq:g-tki}.

\begin{align}
    \text{Gap}_i &= \bar{P}_i - \bar{I}_i \label{eq:g-gap} \\
    \text{TKI}_i &= \frac{\bar{P}_i}{\bar{I}_i} \times 100\% \label{eq:g-tki}
\end{align}

Beberapa bentuk yang sering dibutuhkan: pecahan $\frac{a}{b}$, akar
$\sqrt[3]{x}$, jumlah $\sum_{i=1}^{n}$, hasil kali $\prod_{i=1}^{n}$, integral
$\int_{0}^{1}$, huruf Yunani $\alpha, \beta, \sigma, \mu, \kappa$, serta matriks
berikut.

\begin{equation}
    \label{eq:g-matriks}
    \mathbf{D} =
    \begin{bmatrix}
        d_{11} & d_{12} & \cdots & d_{1n} \\
        d_{21} & d_{22} & \cdots & d_{2n} \\
        \vdots & \vdots & \ddots & \vdots \\
        d_{m1} & d_{m2} & \cdots & d_{mn}
    \end{bmatrix}
\end{equation}

% ==================================================================
\section{Kode Program}
% ==================================================================

\subsection{Kode ditulis langsung}
Kode dapat ditulis langsung di dalam naskah dan dirujuk seperti
\cref{lst:g-python}.

\begin{lstlisting}[language=Python, caption={Contoh kode Python}, label=lst:g-python]
def hitung_rerata(nilai):
    """Menghitung rata-rata dari sebuah senarai angka."""
    if not nilai:
        return 0
    return sum(nilai) / len(nilai)
\end{lstlisting}

\begin{lstlisting}[language=SQL, caption={Contoh kueri SQL}, label=lst:g-sql]
SELECT d.nama       AS dimensi,
       AVG(j.kinerja)     AS rerata_kinerja,
       AVG(j.kepentingan) AS rerata_penting
FROM   jawaban j
JOIN   butir b   ON b.id = j.butir_id
JOIN   dimensi d ON d.id = b.dimensi_id
GROUP  BY d.nama
ORDER  BY rerata_kinerja ASC;
\end{lstlisting}

\begin{lstlisting}[language=Java, caption={Contoh kode Java}, label=lst:g-java]
public class Responden {
    private final String kode;
    private final int[] jawaban;

    public Responden(String kode, int[] jawaban) {
        this.kode = kode;
        this.jawaban = jawaban;
    }

    public double rerata() {
        return Arrays.stream(jawaban)
                     .average()
                     .orElse(0.0);
    }
}
\end{lstlisting}

\begin{lstlisting}[language=PHP, caption={Contoh kode PHP}, label=lst:g-php]
<?php
function hitungGap(array $kin, array $pen): array {
    $gap = [];
    foreach ($kin as $b => $nilai) {
        $gap[$b] = round($nilai - $pen[$b], 2);
    }
    return $gap;
}
\end{lstlisting}

\subsection{Kode dibaca dari berkas}
Cara ini lebih dianjurkan untuk lampiran kode karena isinya selalu sinkron
dengan berkas sumber, dapat dicari, dan tidak membengkakkan ukuran berkas
seperti halnya tangkapan layar. Lihat \cref{lst:g-berkas}.

\lstinputlisting[language=Python, caption={Kode yang dibaca langsung dari berkas \texttt{kode/hitung-sus.py}}, label=lst:g-berkas]{kode/hitung-sus.py}

Sebagian baris saja juga dapat ditampilkan dengan parameter
\texttt{firstline} dan \texttt{lastline}.

\lstinputlisting[language=Python, firstline=7, lastline=18, caption={Cuplikan baris 7--18 dari berkas yang sama}, label=lst:g-cuplikan]{kode/hitung-sus.py}

\subsection{Pseudocode}
Untuk menuliskan algoritma, gunakan \texttt{algorithm2e}. Bentuk ini lebih
disukai penguji daripada kode nyata karena tidak terikat bahasa pemrograman
tertentu. Lihat \cref{alg:g-ipa}.

\begin{algorithm}[H]
\caption{Pemetaan kuadran \textit{Importance-Performance Analysis}}
\label{alg:g-ipa}
\KwMasukan{Senarai butir $B$, skor kepentingan $I$, skor kinerja $P$}
\KwKeluaran{Kuadran setiap butir}
$\bar{I} \leftarrow$ rata-rata seluruh nilai $I$\;
$\bar{P} \leftarrow$ rata-rata seluruh nilai $P$\;
\ForEach{butir $b \in B$}{
    \uIf{$I_b \geq \bar{I}$ \textbf{dan} $P_b < \bar{P}$}{
        $K_b \leftarrow$ ``Prioritas utama''\;
    }
    \uElseIf{$I_b \geq \bar{I}$}{
        $K_b \leftarrow$ ``Pertahankan prestasi''\;
    }
    \uElseIf{$P_b < \bar{P}$}{
        $K_b \leftarrow$ ``Prioritas rendah''\;
    }
    \Else{
        $K_b \leftarrow$ ``Berlebihan''\;
    }
}
\Return $K$\;
\end{algorithm}

% ==================================================================
\section{Sitasi dan Daftar Pustaka}
% ==================================================================

Template ini memakai gaya APA nama-tahun. Gaya diatur melalui perintah
\texttt{\textbackslash gayasitasi} pada berkas \texttt{a0-identitas.tex}.

\subsection{Bentuk sitasi}

\begin{daftarpoin}
    \item \texttt{\textbackslash citep\{kunci\}} menghasilkan sitasi dalam
          kurung: \citep{wijaya2022}. Dipakai bila rujukan bukan bagian dari
          kalimat.
    \item \texttt{\textbackslash citet\{kunci\}} menghasilkan sitasi naratif:
          \citet{wijaya2022}. Dipakai bila nama penulis menjadi subjek kalimat.
    \item Sitasi jamak ditulis dalam satu perintah:
          \citep{wijaya2022, pratama2023, sari2024}.
    \item Sitasi dengan halaman:
          \citep[hlm.~118]{barnes2002}.
    \item Sitasi dengan keterangan awalan:
          \citep[lihat][]{martilla1977}.
    \item Menyebut tahun saja: \citeyearpar{martilla1977}.
    \item Menyebut nama penulis saja: \citeauthor{barnes2002}.
\end{daftarpoin}

\subsection{Kutipan langsung}
Kutipan pendek ditulis menyatu dengan paragraf dan diapit tanda kutip.
Kutipan langsung yang panjang, yaitu lebih dari 40 kata, ditulis sebagai blok
tersendiri yang menjorok tanpa tanda kutip, seperti berikut.

\begin{quote}
    \noindent Tulis kutipan panjang Anda di sini. Blok kutipan ditulis menjorok
    ke dalam, tanpa tanda kutip, dan tetap diakhiri sitasi yang memuat nomor
    halaman sumbernya. Kutipan langsung sebaiknya digunakan seperlunya saja;
    utamakan parafrasa dengan bahasa Anda sendiri agar naskah tidak menjadi
    rangkaian kutipan. \citep[hlm.~120]{barnes2002}
\end{quote}

\subsection{Tipe entri pada berkas .bib}
Berkas \texttt{a7-pustaka.bib} memuat katalog seluruh tipe entri yang lazim
dipakai, yaitu \texttt{@article} untuk artikel jurnal, \texttt{@book} untuk
buku, \texttt{@incollection} untuk bab dalam bunga rampai,
\texttt{@inproceedings} untuk prosiding, \texttt{@mastersthesis} untuk skripsi
atau tesis, \texttt{@techreport} untuk laporan lembaga, \texttt{@manual} untuk
standar dan dokumentasi, serta \texttt{@misc} untuk sumber daring. Salin tipe
yang dibutuhkan lalu isi dengan data sebenarnya.

% Perintah \nocite berikut memaksa seluruh tipe entri ikut tampil pada
% Daftar Pustaka meskipun tidak disitasi di dalam teks, sehingga Anda
% dapat melihat bentuk tampilannya. Entri ini akan hilang sendiri
% begitu galeri dimatikan.
\nocite{contoh-buku, contoh-babbuku, contoh-prosiding, contoh-skripsi,
        contoh-laporan, contoh-standar, contoh-web}

Bentuk tampilan masing-masing tipe dapat Anda lihat langsung pada Daftar
Pustaka naskah ini.

% ==================================================================
\section{Rujukan Silang}
% ==================================================================

Seluruh rujukan silang sebaiknya memakai \texttt{\textbackslash cref} agar kata
``Gambar'', ``Tabel'', atau ``Persamaan'' muncul otomatis dan penomorannya
selalu sinkron: \cref{fig:g-tunggal}, \cref{tab:g-biasa}, \cref{eq:g-rerata},
\cref{lst:g-python}, dan \cref{alg:g-ipa}.

Beberapa objek sekaligus dapat dirujuk dalam satu perintah dan LaTeX akan
menyusun kata sambungnya sendiri: \cref{tab:g-biasa,tab:g-booktabs,tab:g-gabung}.

Rujukan ke rumusan masalah juga tersedia. Perintah \texttt{\textbackslash
cref\{rq:kualitas\}} menghasilkan \cref{rq:kualitas}, sehingga penomoran RQ di
seluruh naskah menyesuaikan sendiri bila urutannya diubah.
 di laporan.tex
% sebelum naskah dikumpulkan. Berkas ini tidak akan ikut tercetak,
% tetapi tetap tersimpan sebagai rujukan.
%==================================================================

% Galeri disertakan setelah lampiran, sehingga tanpa pengaturan ini
% penomorannya akan ikut melanjutkan huruf lampiran (D, E, dst.).
% Blok berikut membuatnya bernomor "G" agar jelas terpisah.
\clearpage
\setcounter{chapter}{0}
\setcounter{figure}{0}
\setcounter{table}{0}
\setcounter{equation}{0}
\renewcommand{\thechapter}{G}
\renewcommand{\thesection}{G.\arabic{section}}
\renewcommand{\thesubsection}{\thesection.\arabic{subsection}}
\renewcommand{\thefigure}{G.\arabic{figure}}
\renewcommand{\thetable}{G.\arabic{table}}
\renewcommand{\theequation}{G.\arabic{equation}}

\chapter*{GALERI ELEMEN (HAPUS SEBELUM DIKUMPULKAN)}
\addcontentsline{toc}{chapter}{GALERI ELEMEN (HAPUS SEBELUM DIKUMPULKAN)}

Bab ini bukan bagian dari naskah Tugas Akhir. Isinya adalah contoh siap salin
untuk seluruh elemen yang lazim dibutuhkan. Matikan dengan memberi tanda
\texttt{\%} pada baris \texttt{\textbackslash include\{zz-galeri-elemen\}} di
berkas \texttt{laporan.tex}.

% ==================================================================
\section{Tabel}
% ==================================================================

\subsection{Tabel biasa}
Keterangan tabel diletakkan di \textbf{atas} tabel. Setiap tabel wajib diberi
\texttt{\textbackslash label} dan disebut dalam teks, misalnya
\cref{tab:g-biasa}.

\begin{table}[H]
    \centering
    \caption{Contoh tabel biasa}
    \label{tab:g-biasa}
    \begin{tabular}{|c|l|c|}
        \hline
        \textbf{No} & \textbf{Keterangan} & \textbf{Nilai} \\ \hline
        1 & Baris pertama & 3,12 \\ \hline
        2 & Baris kedua   & 3,68 \\ \hline
        3 & Baris ketiga  & 3,54 \\ \hline
    \end{tabular}
\end{table}

\subsection{Tabel dengan garis gaya ilmiah (booktabs)}
Sebagian program studi menghendaki tabel tanpa garis vertikal.

\begin{table}[H]
    \centering
    \caption{Contoh tabel gaya booktabs}
    \label{tab:g-booktabs}
    \begin{tabular}{lcc}
        \toprule
        \textbf{Dimensi} & \textbf{Kepentingan} & \textbf{Kinerja} \\
        \midrule
        \textit{Usability}           & 4,41 & 3,12 \\
        \textit{Information quality} & 4,22 & 3,68 \\
        \bottomrule
    \end{tabular}
\end{table}

\subsection{Tabel bersambung antarhalaman (longtable)}
Gunakan \texttt{longtable} bila tabel lebih panjang daripada satu halaman.
Baris kepala akan diulang otomatis, dan halaman sambungan diberi keterangan
``Lanjutan'' sesuai kelaziman naskah Indonesia.

\begin{longtable}{|c|p{6.5cm}|c|c|}
    \caption{Contoh tabel bersambung antarhalaman}
    \label{tab:g-longtable} \\
    \hline
    \textbf{No} & \textbf{Pernyataan} & \textbf{Kepentingan} & \textbf{Kinerja} \\ \hline
    \endfirsthead

    \multicolumn{4}{l}{\footnotesize\textit{Lanjutan \cref{tab:g-longtable}}} \\
    \hline
    \textbf{No} & \textbf{Pernyataan} & \textbf{Kepentingan} & \textbf{Kinerja} \\ \hline
    \endhead

    \hline
    \multicolumn{4}{r}{\footnotesize\textit{Bersambung ke halaman berikutnya}} \\
    \endfoot

    \hline
    \endlastfoot

    1  & Saya merasa sistem mudah dipelajari & 4,44 & 3,10 \\ \hline
    2  & Interaksi dengan sistem jelas dan mudah dipahami & 4,38 & 3,21 \\ \hline
    3  & Sistem mudah dioperasikan & 4,41 & 3,15 \\ \hline
    4  & Navigasi antarhalaman terasa mudah & 4,61 & 3,03 \\ \hline
    5  & Tampilan sistem sesuai dengan jenis layanannya & 4,02 & 3,44 \\ \hline
    6  & Tampilan sistem menarik & 3,88 & 3,52 \\ \hline
    7  & Sistem menyampaikan kesan kompeten & 4,12 & 3,40 \\ \hline
    8  & Sistem menciptakan pengalaman yang menyenangkan & 3,95 & 3,33 \\ \hline
    9  & Informasi yang disajikan akurat & 4,30 & 3,71 \\ \hline
    10 & Informasi yang disajikan tepercaya & 4,28 & 3,69 \\ \hline
    11 & Informasi selalu diperbarui & 4,19 & 3,58 \\ \hline
    12 & Informasi relevan dengan kebutuhan & 4,25 & 3,66 \\ \hline
    13 & Informasi mudah dipahami & 4,21 & 3,74 \\ \hline
    14 & Informasi disajikan pada tingkat rincian yang tepat & 4,08 & 3,70 \\ \hline
    15 & Informasi disajikan dalam format yang sesuai & 4,15 & 3,68 \\ \hline
    16 & Sistem memberikan rasa aman & 4,09 & 3,61 \\ \hline
    17 & Data pribadi terasa terlindungi & 4,18 & 3,55 \\ \hline
    18 & Sistem memudahkan komunikasi dengan pengelola & 3,97 & 3,42 \\ \hline
    19 & Sistem menciptakan rasa keterhubungan & 3,82 & 3,39 \\ \hline
    20 & Layanan diberikan sesuai yang dijanjikan & 4,11 & 3,58 \\ \hline
\end{longtable}

\subsection{Tabel dengan sel gabung (multirow dan multicolumn)}

\begin{table}[H]
    \centering
    \caption{Contoh tabel dengan sel gabung}
    \label{tab:g-gabung}
    \begin{tabular}{|c|c|c|c|c|}
        \hline
        \multirow{2}{*}{\textbf{Butir}} & \multicolumn{2}{c|}{\textbf{Sebelum}} & \multicolumn{2}{c|}{\textbf{Sesudah}} \\ \cline{2-5}
         & \textbf{Rerata} & \textbf{SB} & \textbf{Rerata} & \textbf{SB} \\ \hline
        B1 & 3,10 & 0,82 & 4,05 & 0,61 \\ \hline
        B2 & 3,22 & 0,77 & 4,12 & 0,58 \\ \hline
    \end{tabular}
\end{table}

\subsection{Tabel berkolom otomatis (tabularx)}
Lebar kolom bertipe \texttt{X} menyesuaikan sendiri agar tabel tepat memenuhi
lebar teks.

\begin{table}[H]
    \centering
    \caption{Contoh tabel dengan lebar kolom otomatis}
    \label{tab:g-tabularx}
    \begin{tabularx}{\textwidth}{|c|X|X|}
        \hline
        \textbf{No} & \textbf{Temuan} & \textbf{Tindak lanjut} \\ \hline
        1 & Menu utama tidak konsisten antarhalaman sehingga pengguna kehilangan orientasi. & Menyeragamkan struktur menu pada seluruh halaman. \\ \hline
        2 & Pesan kesalahan tidak menjelaskan penyebab masalah. & Menulis ulang pesan kesalahan dengan bahasa yang spesifik. \\ \hline
    \end{tabularx}
\end{table}

\subsection{Tabel dengan catatan kaki (threeparttable)}
Gunakan bila tabel memerlukan keterangan simbol atau sumber.

\begin{table}[H]
    \centering
    \begin{threeparttable}
        \caption{Contoh tabel dengan catatan kaki}
        \label{tab:g-catatan}
        \begin{tabular}{|l|c|c|}
            \hline
            \textbf{Dimensi} & \textbf{Rerata} & \textbf{Kategori} \\ \hline
            \textit{Usability}\tnote{a}           & 3,12 & Cukup \\ \hline
            \textit{Information quality}          & 3,68 & Baik  \\ \hline
            \textit{Service interaction}\tnote{b} & 3,54 & Baik  \\ \hline
        \end{tabular}
        \begin{tablenotes}\footnotesize
            \item[a] Terdiri atas 8 butir pernyataan.
            \item[b] Terdiri atas 7 butir pernyataan.
            \item Sumber: Hasil pengolahan data penelitian (2026).
        \end{tablenotes}
    \end{threeparttable}
\end{table}

\subsection{Tabel yang terlalu lebar}
Bila tabel sedikit melebihi lebar teks, perkecil jarak antarkolom dengan
\texttt{\textbackslash tabcolsep} atau bungkus dengan \texttt{adjustbox}.

\begin{table}[H]
    \centering
    \caption{Contoh tabel lebar yang dikecilkan otomatis}
    \label{tab:g-lebar}
    \footnotesize
    \begin{adjustbox}{max width=\textwidth}
    \begin{tabular}{|c|c|c|c|c|c|c|c|c|c|c|}
        \hline
        \textbf{Butir} & \textbf{R1} & \textbf{R2} & \textbf{R3} & \textbf{R4} & \textbf{R5} & \textbf{R6} & \textbf{R7} & \textbf{R8} & \textbf{Rerata} & \textbf{Simpangan Baku} \\ \hline
        B1 & 4 & 5 & 4 & 3 & 4 & 5 & 4 & 4 & 4,13 & 0,64 \\ \hline
        B2 & 3 & 4 & 4 & 4 & 5 & 4 & 3 & 4 & 3,88 & 0,64 \\ \hline
    \end{tabular}
    \end{adjustbox}
\end{table}

\subsection{Tabel lanskap}
Paket \texttt{pdflscape} ikut memutar halaman pada berkas PDF, sehingga pembaca
tidak perlu memiringkan kepala.

\begin{landscape}
\begin{table}[H]
    \centering
    \caption{Contoh tabel pada halaman lanskap}
    \label{tab:g-lanskap}
    \footnotesize
    \begin{tabular}{|c|p{4cm}|p{4cm}|p{4cm}|p{4cm}|}
        \hline
        \textbf{No} & \textbf{Kondisi berjalan} & \textbf{Masalah} & \textbf{Usulan} & \textbf{Dampak yang diharapkan} \\ \hline
        1 & Menu tersebar pada tiga tingkat & Pengguna kesulitan menemukan menu & Menyusun ulang menjadi dua tingkat & Waktu pencarian menu berkurang \\ \hline
        2 & Alur pengisian tidak menampilkan tahapan & Pengguna tidak tahu posisinya & Menambahkan penanda tahapan & Kesalahan pengisian berkurang \\ \hline
    \end{tabular}
\end{table}
\end{landscape}

% ==================================================================
\section{Gambar}
% ==================================================================

\subsection{Gambar tunggal}
Keterangan gambar diletakkan di \textbf{bawah} gambar. Lihat
\cref{fig:g-tunggal}.

\begin{figure}[H]
    \centering
    \includegraphics[width=0.6\linewidth]{contoh-diagram}
    \caption{Contoh gambar tunggal}
    \label{fig:g-tunggal}
\end{figure}

\subsection{Gambar dengan keterangan sumber}
Gunakan perintah \texttt{\textbackslash sumber} bila gambar dikutip dari
sumber lain.

\begin{figure}[H]
    \centering
    \includegraphics[width=0.55\linewidth]{contoh-diagram}
    \caption{Contoh gambar yang dikutip dari sumber lain}
    \label{fig:g-sumber}
\end{figure}
\sumber{\citet{barnes2002}}

\subsection{Dua gambar berdampingan}
Gunakan \texttt{subcaption} agar tiap bagian memperoleh label (a) dan (b) yang
dapat dirujuk terpisah, misalnya \cref{fig:g-kiri} dan \cref{fig:g-kanan}.

\begin{figure}[H]
    \centering
    \begin{subfigure}[t]{0.47\linewidth}
        \centering
        \includegraphics[width=\linewidth]{contoh-tangkapan-a}
        \caption{Tampilan sebelum perbaikan}
        \label{fig:g-kiri}
    \end{subfigure}
    \hfill
    \begin{subfigure}[t]{0.47\linewidth}
        \centering
        \includegraphics[width=\linewidth]{contoh-tangkapan-b}
        \caption{Tampilan sesudah perbaikan}
        \label{fig:g-kanan}
    \end{subfigure}
    \caption{Perbandingan tampilan antarmuka}
    \label{fig:g-berdampingan}
\end{figure}

\subsection{Gambar tinggi satu halaman}
Untuk gambar yang tinggi, seperti diagram alir panjang, gunakan
\texttt{height} agar tidak meluber ke halaman berikutnya.

\begin{figure}[H]
    \centering
    \includegraphics[height=0.78\textheight,keepaspectratio]{contoh-tinggi}
    \caption{Contoh gambar tinggi satu halaman}
    \label{fig:g-tinggi}
\end{figure}

\subsection{Gambar lanskap}

\begin{landscape}
\begin{figure}[H]
    \centering
    \includegraphics[width=0.92\linewidth]{contoh-lanskap}
    \caption{Contoh gambar pada halaman lanskap}
    \label{fig:g-lanskap}
\end{figure}
\end{landscape}

\subsection{Catatan penempatan gambar}
Parameter \texttt{[H]} memaksa gambar muncul tepat di posisi penulisannya.
Parameter \texttt{[htbp]} memberi LaTeX kebebasan menempatkan gambar pada
posisi terbaik, sehingga halaman tidak menyisakan ruang kosong. Gunakan
\texttt{[H]} bila urutan gambar penting, dan \texttt{[htbp]} bila
mengutamakan kerapian halaman.

% ==================================================================
\section{Persamaan}
% ==================================================================

Persamaan dalam baris teks ditulis seperti $\bar{x} = \sum x_i / n$, sedangkan
persamaan bernomor ditulis terpisah seperti \cref{eq:g-rerata}.

\begin{equation}
    \label{eq:g-rerata}
    \bar{x} = \frac{1}{n}\sum_{i=1}^{n} x_{i}
\end{equation}

Persamaan bertingkat ditulis dengan lingkungan \texttt{align}, dan setiap baris
memperoleh nomor tersendiri seperti \cref{eq:g-gap} dan \cref{eq:g-tki}.

\begin{align}
    \text{Gap}_i &= \bar{P}_i - \bar{I}_i \label{eq:g-gap} \\
    \text{TKI}_i &= \frac{\bar{P}_i}{\bar{I}_i} \times 100\% \label{eq:g-tki}
\end{align}

Beberapa bentuk yang sering dibutuhkan: pecahan $\frac{a}{b}$, akar
$\sqrt[3]{x}$, jumlah $\sum_{i=1}^{n}$, hasil kali $\prod_{i=1}^{n}$, integral
$\int_{0}^{1}$, huruf Yunani $\alpha, \beta, \sigma, \mu, \kappa$, serta matriks
berikut.

\begin{equation}
    \label{eq:g-matriks}
    \mathbf{D} =
    \begin{bmatrix}
        d_{11} & d_{12} & \cdots & d_{1n} \\
        d_{21} & d_{22} & \cdots & d_{2n} \\
        \vdots & \vdots & \ddots & \vdots \\
        d_{m1} & d_{m2} & \cdots & d_{mn}
    \end{bmatrix}
\end{equation}

% ==================================================================
\section{Kode Program}
% ==================================================================

\subsection{Kode ditulis langsung}
Kode dapat ditulis langsung di dalam naskah dan dirujuk seperti
\cref{lst:g-python}.

\begin{lstlisting}[language=Python, caption={Contoh kode Python}, label=lst:g-python]
def hitung_rerata(nilai):
    """Menghitung rata-rata dari sebuah senarai angka."""
    if not nilai:
        return 0
    return sum(nilai) / len(nilai)
\end{lstlisting}

\begin{lstlisting}[language=SQL, caption={Contoh kueri SQL}, label=lst:g-sql]
SELECT d.nama       AS dimensi,
       AVG(j.kinerja)     AS rerata_kinerja,
       AVG(j.kepentingan) AS rerata_penting
FROM   jawaban j
JOIN   butir b   ON b.id = j.butir_id
JOIN   dimensi d ON d.id = b.dimensi_id
GROUP  BY d.nama
ORDER  BY rerata_kinerja ASC;
\end{lstlisting}

\begin{lstlisting}[language=Java, caption={Contoh kode Java}, label=lst:g-java]
public class Responden {
    private final String kode;
    private final int[] jawaban;

    public Responden(String kode, int[] jawaban) {
        this.kode = kode;
        this.jawaban = jawaban;
    }

    public double rerata() {
        return Arrays.stream(jawaban)
                     .average()
                     .orElse(0.0);
    }
}
\end{lstlisting}

\begin{lstlisting}[language=PHP, caption={Contoh kode PHP}, label=lst:g-php]
<?php
function hitungGap(array $kin, array $pen): array {
    $gap = [];
    foreach ($kin as $b => $nilai) {
        $gap[$b] = round($nilai - $pen[$b], 2);
    }
    return $gap;
}
\end{lstlisting}

\subsection{Kode dibaca dari berkas}
Cara ini lebih dianjurkan untuk lampiran kode karena isinya selalu sinkron
dengan berkas sumber, dapat dicari, dan tidak membengkakkan ukuran berkas
seperti halnya tangkapan layar. Lihat \cref{lst:g-berkas}.

\lstinputlisting[language=Python, caption={Kode yang dibaca langsung dari berkas \texttt{kode/hitung-sus.py}}, label=lst:g-berkas]{kode/hitung-sus.py}

Sebagian baris saja juga dapat ditampilkan dengan parameter
\texttt{firstline} dan \texttt{lastline}.

\lstinputlisting[language=Python, firstline=7, lastline=18, caption={Cuplikan baris 7--18 dari berkas yang sama}, label=lst:g-cuplikan]{kode/hitung-sus.py}

\subsection{Pseudocode}
Untuk menuliskan algoritma, gunakan \texttt{algorithm2e}. Bentuk ini lebih
disukai penguji daripada kode nyata karena tidak terikat bahasa pemrograman
tertentu. Lihat \cref{alg:g-ipa}.

\begin{algorithm}[H]
\caption{Pemetaan kuadran \textit{Importance-Performance Analysis}}
\label{alg:g-ipa}
\KwMasukan{Senarai butir $B$, skor kepentingan $I$, skor kinerja $P$}
\KwKeluaran{Kuadran setiap butir}
$\bar{I} \leftarrow$ rata-rata seluruh nilai $I$\;
$\bar{P} \leftarrow$ rata-rata seluruh nilai $P$\;
\ForEach{butir $b \in B$}{
    \uIf{$I_b \geq \bar{I}$ \textbf{dan} $P_b < \bar{P}$}{
        $K_b \leftarrow$ ``Prioritas utama''\;
    }
    \uElseIf{$I_b \geq \bar{I}$}{
        $K_b \leftarrow$ ``Pertahankan prestasi''\;
    }
    \uElseIf{$P_b < \bar{P}$}{
        $K_b \leftarrow$ ``Prioritas rendah''\;
    }
    \Else{
        $K_b \leftarrow$ ``Berlebihan''\;
    }
}
\Return $K$\;
\end{algorithm}

% ==================================================================
\section{Sitasi dan Daftar Pustaka}
% ==================================================================

Template ini memakai gaya APA nama-tahun. Gaya diatur melalui perintah
\texttt{\textbackslash gayasitasi} pada berkas \texttt{a0-identitas.tex}.

\subsection{Bentuk sitasi}

\begin{daftarpoin}
    \item \texttt{\textbackslash citep\{kunci\}} menghasilkan sitasi dalam
          kurung: \citep{wijaya2022}. Dipakai bila rujukan bukan bagian dari
          kalimat.
    \item \texttt{\textbackslash citet\{kunci\}} menghasilkan sitasi naratif:
          \citet{wijaya2022}. Dipakai bila nama penulis menjadi subjek kalimat.
    \item Sitasi jamak ditulis dalam satu perintah:
          \citep{wijaya2022, pratama2023, sari2024}.
    \item Sitasi dengan halaman:
          \citep[hlm.~118]{barnes2002}.
    \item Sitasi dengan keterangan awalan:
          \citep[lihat][]{martilla1977}.
    \item Menyebut tahun saja: \citeyearpar{martilla1977}.
    \item Menyebut nama penulis saja: \citeauthor{barnes2002}.
\end{daftarpoin}

\subsection{Kutipan langsung}
Kutipan pendek ditulis menyatu dengan paragraf dan diapit tanda kutip.
Kutipan langsung yang panjang, yaitu lebih dari 40 kata, ditulis sebagai blok
tersendiri yang menjorok tanpa tanda kutip, seperti berikut.

\begin{quote}
    \noindent Tulis kutipan panjang Anda di sini. Blok kutipan ditulis menjorok
    ke dalam, tanpa tanda kutip, dan tetap diakhiri sitasi yang memuat nomor
    halaman sumbernya. Kutipan langsung sebaiknya digunakan seperlunya saja;
    utamakan parafrasa dengan bahasa Anda sendiri agar naskah tidak menjadi
    rangkaian kutipan. \citep[hlm.~120]{barnes2002}
\end{quote}

\subsection{Tipe entri pada berkas .bib}
Berkas \texttt{a7-pustaka.bib} memuat katalog seluruh tipe entri yang lazim
dipakai, yaitu \texttt{@article} untuk artikel jurnal, \texttt{@book} untuk
buku, \texttt{@incollection} untuk bab dalam bunga rampai,
\texttt{@inproceedings} untuk prosiding, \texttt{@mastersthesis} untuk skripsi
atau tesis, \texttt{@techreport} untuk laporan lembaga, \texttt{@manual} untuk
standar dan dokumentasi, serta \texttt{@misc} untuk sumber daring. Salin tipe
yang dibutuhkan lalu isi dengan data sebenarnya.

% Perintah \nocite berikut memaksa seluruh tipe entri ikut tampil pada
% Daftar Pustaka meskipun tidak disitasi di dalam teks, sehingga Anda
% dapat melihat bentuk tampilannya. Entri ini akan hilang sendiri
% begitu galeri dimatikan.
\nocite{contoh-buku, contoh-babbuku, contoh-prosiding, contoh-skripsi,
        contoh-laporan, contoh-standar, contoh-web}

Bentuk tampilan masing-masing tipe dapat Anda lihat langsung pada Daftar
Pustaka naskah ini.

% ==================================================================
\section{Rujukan Silang}
% ==================================================================

Seluruh rujukan silang sebaiknya memakai \texttt{\textbackslash cref} agar kata
``Gambar'', ``Tabel'', atau ``Persamaan'' muncul otomatis dan penomorannya
selalu sinkron: \cref{fig:g-tunggal}, \cref{tab:g-biasa}, \cref{eq:g-rerata},
\cref{lst:g-python}, dan \cref{alg:g-ipa}.

Beberapa objek sekaligus dapat dirujuk dalam satu perintah dan LaTeX akan
menyusun kata sambungnya sendiri: \cref{tab:g-biasa,tab:g-booktabs,tab:g-gabung}.

Rujukan ke rumusan masalah juga tersedia. Perintah \texttt{\textbackslash
cref\{rq:kualitas\}} menghasilkan \cref{rq:kualitas}, sehingga penomoran RQ di
seluruh naskah menyesuaikan sendiri bila urutannya diubah.
